\section{Aislamiento del circuito frío frente al exterior húmedo (INT-34)}
\label{sec:sd4-anticondensacion-int34}
Se separa la condición interior limpia de ENV-25 de la exposición exterior de los tramos fríos. El cálculo con cámara a 24 grados C no acredita una conexión de agua, serpentín o conducto expuesta a 35 grados C y 80\,\% HR. Se fija para obra una superficie fría mínima de cálculo de 5 grados C, igual al agua de suministro proyectada, y conductividad máxima de aislamiento seco de 0,040 W/(m K). Estas condiciones no son temperatura superficial medida ni propiedades de un material ya comprado. El cálculo de rocío mediante aproximación $g=\ln(0,80)+17,62(35)/(243,12+35)$, $T_r=243,12g/(17,62-g)$ da 31,0283 grados C. Se exige margen de superficie de al menos 1 K sobre ese rocío en campo, antes de resolver puentes.

Para la comprobación conservadora plana se desprecia resistencia interior, lo que aumenta la transferencia hacia la superficie fría, y se adopta coeficiente exterior mínimo propio de 4 W/(m$^2$K), con sensibilidad a 8. Se obtiene $R=d/0,040+1/h_o$, $q''=(35-5)/R$ y $T_{so}=35-q''/h_o$. Con 25 mm y $h_o=4$, la superficie resulta 26,4286 grados C; con 50 mm, 30,0000 grados C: ambas están por debajo del rocío de proyecto. El valor de 25 mm de ENV-25 se mantiene sólo para su entorno interior tratado y sus condiciones verificadas; no se extiende al exterior. Los 50 mm publicados para una familia CFK tampoco acreditan el punto húmedo de una UTA distinta.

Se selecciona para los tramos fríos expuestos al exterior una reserva de obra de aislamiento no combustible de al menos 100 mm y $\lambda\le0,040$ W/(m K), con barrera de vapor continua en cara caliente, revestimiento resistente al ambiente marino y soportes con rotura térmica. A $h_o=4$, $R=2,75$ m$^2$K/W, $q''=10,9091$ W/m$^2$ y $T_{so}=32,2727$ grados C, margen 1,2444 K. A $h_o=8$ resulta 33,5714 grados C. El espesor mínimo teórico para margen de 1 K bajo estos supuestos es aproximadamente 91,0 mm; se fija 100 mm de obra, sin convertirlo en espesor certificado de una manta o clase al fuego adquirida. Si la conductividad real húmeda, el revestimiento o el coeficiente local invalidan la cota, se aumenta o cambia el conjunto antes de ejecutar.

En DN315 se reserva diámetro exterior de 515 mm con 100 mm por lado, antes de revestimiento, herrajes y tolerancias. Esa dimensión se concilia con las reservas de UTA/plenum y bases vigentes; no se atribuye cabida a todo el serpentín por el diámetro de un tubo. El aislamiento exterior nuevo se limita al circuito frío antes de recalentamiento y a tubería, válvulas y bandejas frías: no cubre calefactor, tapas eléctricas ni espacios de ventilación prohibidos por sus manuales. Al pasar a espacio tratado se distingue cada detalle y no se elimina barrera de vapor en una junta oculta. La clasificación del aire de banco y la prevención de retorno H$_2$ se resuelven por sus dispositivos, sin atribuirlas al revestimiento.

Los cálculos de campo no acreditan bridas, vástagos, desagües, fijaciones, bandeja ni puentes metálicos. Cada detalle se diseña con continuidad de vapor y separación del ambiente húmedo, sin descarga de condensado sobre tableros ni comunicación gaseosa por drenaje. Se cuantifican materiales y longitudes con el trazado final, sin comprar por longitud equivalente hidráulica. La ganancia térmica de las superficies reales entra en el serpentín y la fuente de frío una sola vez: 22,8604 kW por banco corresponde al punto previo de rocío 8,5 grados C y no acredita el nuevo presupuesto instrumental ni absorbe esta ganancia desconocida. Como contraste ideal a salida saturada de 7,3925 grados C y 900 m$^3$/h, la demanda de frío resulta 23,516928 kW por banco, agua a salto 5 K de 4,041903 m$^3$/h por ruta y condensado 22,203622 kg/h por banco. La cota de lectura instrumental no se convierte en temperatura real de salida adoptada: faltan margen de instalación/presión, curva húmeda, bypass, ensuciamiento y ganancia del trazado. Este contraste tampoco es potencia eléctrica de una fuente de frío seleccionada. La recepción prevista incluye inspección de vapor/soportes, temperaturas superficiales, ensayo de drenaje y alarma de agua en el estado húmedo de diseño. La especificación resuelve la frontera anticondensación de campo; materiales comerciales, puentes, capacidad húmeda y carga eléctrica completa siguen pendientes.
